\documentclass{article}


\begin{document}

\begin{center}
{\bf Slopes in submodules of a free module.}
\end{center}

\begin{center}
Jose Mar\'{\i}a Ucha Enr\'{\i}quez.

Dpto. Algebra. Universidad de Sevilla.

ucha@algebra.us.es

954 55 69 69, FAX: 954 55 69 38
\end{center}


\noindent
Let ${\cal A} = A_n ({\bf C}) = {\bf C}[X]<D>= {\bf C}[x_1, ..., x_n]
<\partial_1, ..., \partial_n>$ 
be the $n$-th Weyl algebra, and $L: {\bf Q}^2 \longrightarrow 
{\bf Q}$ a linear form with possitive coeficients.
  
Let us consider the free module ${\cal A}^p$ and ${\cal N}$ a submodule.
The $L$-order of an element is the maximum of the values of 
$L(|\beta|,\beta_1 - \alpha_1)$ over the exponents $(\alpha, \beta)$ in any 
component of the element.  This order induces a $L$-filtration and an 
associated graduated module $gr^L ({\cal A}^p)$.
    
We will say that $L$ is a slope of ${\cal M} = {\cal A}^p/{\cal N}$
if the ideal
$$
\sqrt{A^{\cal M}_L} = \sqrt{Ann_{gr^L ({\cal A})} (gr^L ({\cal A}^p)
/gr^L({\cal N}))}
$$
is not $(F,V)$-bihomogeneous, where $F$ 
is the graduation respect to the order on $D$
and $V$ is the graduation of Malgrange-Kashiwara.

In this work we generalize ``How to calculate the slope of a ${\cal D}$-
module'' from Assi-Castro-Granger, and  
``Homogenising differential operators'' from Castro-Narvaez. 
We provide an algorithm of effective calculus of 
slopes as well, that we have implemented in COMMON LISP. 
Eventually, we use MACAULAY to calculate radicals of ideals of 
polynomial rings. 



Some concrete examples are treated:
\begin{itemize} 
\item Syzygy modules. 
\item Modules of the form 
${\cal O}[1/f]/{\cal O}$ with $f$ a surface.
\item ${\cal A}^2 /{\cal N}$ with ${\cal N} 
= I_1 \oplus I_2$, and each $I_i$ a left ideal of ${\cal A}$ with known 
slopes.
\end{itemize} 
   


\end{document}








